\DocumentMetadata{
  pdfversion=2.0,
  pdfstandard=ua-2,
  lang=en-US,
  tagging=on,
  tagging-setup={math/setup=mathml-SE,tabsorder=structure}
}
\documentclass[11pt,letterpaper]{article}
\usepackage[margin=0.68in,includefoot]{geometry}
\usepackage{newcomputermodern}
\usepackage{amsmath,amscd,mathtools}
\usepackage{tikz,adjustbox}
\providecommand{\square}{\mdlgwhtsquare}
\usepackage[hidelinks]{hyperref}
\hypersetup{
  pdftitle={Analysis First-Year Exam, May 2018, Part 2},
  pdfauthor={University of Florida Department of Mathematics},
  pdfsubject={Graduate mathematics examination},
  pdfdisplaydoctitle=true,
  bookmarksopen=true
}
\setcounter{secnumdepth}{0}
\setlength{\parindent}{0pt}
\setlength{\parskip}{0.28em}
\setlength{\emergencystretch}{3em}
\setlength{\leftmargini}{2.15em}
\setlength{\leftmarginii}{2.65em}
\setlength{\leftmarginiii}{3em}
\pagestyle{empty}
\raggedbottom
\allowdisplaybreaks
\definecolor{ProblemConcern}{HTML}{7A1F1F}
\newenvironment{problemconcern}{%
  \par\tagstructbegin{tag=Aside}\begingroup\color{ProblemConcern}
}{%
  \par\endgroup\tagstructend\nopagebreak[4]\vspace{0.12em}
}
\NewTaggingSocketPlug{block/list/label}{examproblem}{%
  \tagstructbegin{tag=Lbl}\tagstructend%
  \tagstructbegin{tag=\LBody}%
  \tagstructbegin{tag=\ProblemHeadingTag,title-o={\ProblemHeadingText}}%
  \tagmcbegin{tag=\ProblemHeadingTag,actualtext={\ProblemHeadingText}}%
  #2%
  \tagmcend%
  \tagstructend%
}
\newenvironment{examproblems}[2]{%
  \def\ProblemHeadingTag{#2}%
  \AssignTaggingSocketPlug{block/list/label}{examproblem}%
  \begin{enumerate}%
  \setcounter{enumi}{#1}%
  \renewcommand{\labelenumi}{\textbf{\theenumi.}}%
  \setlength{\topsep}{0.35em}%
  \setlength{\partopsep}{0pt}%
  \setlength{\itemsep}{0.48em}%
  \setlength{\parsep}{0.18em}%
}{\end{enumerate}}
\newenvironment{examsubparts}{%
  \AssignTaggingSocketPlug{block/list/label}{default}%
  \begin{enumerate}%
  \setlength{\topsep}{0.18em}%
  \setlength{\partopsep}{0pt}%
  \setlength{\itemsep}{0.16em}%
  \setlength{\parsep}{0.1em}%
}{\end{enumerate}}
\newenvironment{examinstructions}{%
  \par\begingroup\itshape
}{%
  \par\endgroup\vspace{0.18em}
}
\makeatletter
\renewcommand\subsection{\@startsection{subsection}{2}{0pt}%
  {-1.25ex plus -.25ex minus -.1ex}%
  {0.55ex plus .1ex}%
  {\normalfont\large\bfseries}}
\renewcommand{\@maketitle}{%
  \newpage\null\vskip -1em%
  {\footnotesize\bfseries
    \href{https://www.ufl.edu/}{University of Florida}, \href{https://math.ufl.edu/}{Department of Mathematics}\par}%
  \vskip -0.5em%
  \tagstructbegin{tag=H1,title={Analysis First-Year Exam, May 2018, Part 2}}%
  \tagpdfparaOff%
  \tagmcbegin{tag=H1}%
    {\LARGE\bfseries\href{https://gma.math.ufl.edu/exams/analysis-first-year/}{Analysis First-Year Exam}, May 2018, Part 2\par}%
  \tagmcend%
  \tagpdfparaOn%
  \tagstructend%
  \vskip 0.25em%
  \tagmcbegin{artifact}%
    \hrule height 1pt%
  \tagmcend%
  \vskip 1em%
}
\makeatother
\title{Analysis First-Year Exam, May 2018, Part 2}
\author{}
\date{}
\begin{document}
\maketitle
\thispagestyle{empty}
\begin{examinstructions}
Answer FOUR questions in detail. State carefully any results used without proof.
\end{examinstructions}
\begin{examproblems}{0}{H2}
\def\ProblemHeadingText{Problem 1.}
\item
This problem has two parts.
\begin{examsubparts}
\item[\textnormal{(i)}]
Let \(g_n \to g\) and \(h_n \to h\) uniformly on~\(X\). Prove that if~\(g\) and~\(h\) are bounded, then \(g_nh_n \to gh\) uniformly on~\(X\).
\item[\textnormal{(ii)}]
Let \(f:[a,b]\to\mathbb{R}\) be continuous. By considering the square-root or otherwise, prove that if~\(f\) is non-negative then there exists a sequence of non-negative polynomials converging uniformly to~\(f\) on \([a,b]\).
\end{examsubparts}
\def\ProblemHeadingText{Problem 2.}
\item
Let~\(\mathcal{F}\) be an equicontinuous family of real-valued functions on the connected space~\(X\). By considering the set of points at which the family~\(\mathcal{F}\) is bounded, prove that if~\(\mathcal{F}\) is bounded at one point of~\(X\) then it is bounded at all points of~\(X\).
\def\ProblemHeadingText{Problem 3.}
\item
\((f_n:n\geq 0)\) is a sequence of measurable real-valued functions. Prove that each of the following sets is measurable:
\begin{examsubparts}
\item[\textnormal{(a)}]
the set~\(A\) comprising all points~\(\omega\) for which the sequence of values \(f_n(\omega)\) is eventually positive;
\item[\textnormal{(b)}]
the set~\(B\) comprising all points~\(\omega\) for which the sequence of values \(f_n(\omega)\) changes sign infinitely often.
\end{examsubparts}
\def\ProblemHeadingText{Problem 4.}
\item
\begin{problemconcern}
\textbf{Warning: Suspected error.} The source clearly says “Lebesgue-measurable functions,” but measurability and uniform convergence alone do not ensure the stated integrability conclusion, even when the measure is finite. A missing integrability hypothesis may have been intended, but the correction is not uniquely determined, so the source wording has been retained.
\end{problemconcern}
\((\Omega,\mathcal{F},\mu)\) is a measure space on which \((f_n:n\geq 0)\) is a sequence of Lebesgue-measurable functions converging uniformly to~\(f\). Prove that if \(\mu(\Omega)<\infty\) then it follows that~\(f\) is integrable and \(\int_\Omega f_n\,d\mu\to\int_\Omega f\,d\mu\). Show by example that if \(\mu(\Omega)=\infty\) then the same conclusion may not follow.
\def\ProblemHeadingText{Problem 5.}
\item
Let \((f_n:n\geq 0)\) be a uniformly-bounded sequence of Riemann-integrable functions on \([a,b]\). If \(f_n\to 0\) pointwise, does it follow that \(\int_a^b f_n\to 0\)? Proof or counterexample required.
\end{examproblems}
\end{document}
